\documentclass[../../../include/open-logic-section]{subfiles}

\begin{document}

\olfileid{sth}{ordinals}{vn}
\olsection[فون نوئمان دی بناوٹ]{ترتیبی عدداں دی فون نوئمان والی بناوٹ}

\olref[sth][ordinals][iso]{thm:woalwayscomparable} توں اک خیال آؤندا اے۔
اسیں کجھ شےواں متعارف کرا سکدے آں، جیہناں نوں \emph{ترتیب دیاں قسماں}
آکھاں گے، تے جیہڑیاں حُسنِ ترتیب والیاں ساختاں دی تھاں ورتیاں جان۔
جے حُسنِ ترتیب والی ساخت دی ترتیب دی قسم لئی $\ordtype{A, <}$ لکھدے آں،
جیہڑی ساخت $\tuple{A, <}$ اے، تاں سانوں ہِیٹھلے دو اصول ملن دی آس ہووے گی:
\begin{align*}
	\ordtype{A, <} = \ordtype{B, \lessdot} &
	\text{ اُدوں تے صرف اُدوں جدوں } \ordeq{\tuple{A, <}}{\tuple{B, \lessdot}}\\
	\ordtype{A, <} < \ordtype{B, \lessdot}&
	\text{ اُدوں تے صرف اُدوں جدوں }\ordeq{\tuple{A, <}}{\tuple{B_b, \lessdot_b}}\text{ کسے ایہو جیہے }b \in B
\end{align*}
نالے اسیں ترتیب دیاں قسماں نوں \emph{کجھ خاص سیٹاں دے طور اُتے}
متعارف کراؤن دی آس رکھ سکدے آں، جیویں قدرتی عدداں نوں کجھ خاص سیٹاں
دے طور اُتے متعارف کرا سکدے آں۔

ایہ کرن دا سبھ توں عام طریقہ، تے اوہی طریقہ جیہڑا اسیں ورتاں گے،
ترتیب دیاں قسماں نوں کجھ \emph{معیاری} حُسنِ ترتیب والے سیٹاں نال
تعریف کرنا اے۔ ایہ معیاری سیٹ سبھ توں پہلاں فون نوئمان نے متعارف کرائے:

\begin{defn}
سیٹ $A$ \emph{متعدی} اُدوں تے صرف اُدوں اے جدوں
$(\forall x \in A)x \subseteq A$۔ فیر $A$ \emph{ترتیبی عدد} اُدوں تے
صرف اُدوں اے جدوں $A$ متعدی اے تے $\in$ اینوں حُسنِ ترتیب وچ رکھدا اے۔
\end{defn}
\noindent
اگے اسیں ترتیبی عدداں لئی یونانی حرف ورتاں گے۔ تعریف توں جھٹ نکلدا اے
کہ جے $\alpha$ ترتیبی عدد اے، تاں $\tuple{\alpha, \in_\alpha}$
حُسنِ ترتیب والی ساخت اے، جتھے $\in_\alpha =
\Setabs{\tuple{x, y} \in \alpha^2}{x \in y}$۔ سو نشان نوں ذرا کھُلّ
نال ورت کے آکھ سکدے آں کہ $\alpha$ \emph{آپ} حُسنِ ترتیب والی ساخت اے۔

ساڈے پہلے کجھ ترتیبی عدد ایہ نیں:
\[
	\emptyset, \{\emptyset\},
	\{\emptyset, \{\emptyset\}\},
	\{\emptyset, \{\emptyset\}, \{\emptyset, \{\emptyset\}\}\}, \ldots
\]
تسی ویکھو گے کہ ایہ اوہی پہلے کجھ ترتیبی عدد نیں جیہڑے ساڈے لامتناہیت دے
مسلّمے وچ آئے سن، یعنی فون نوئمان دی $\omega$ دی تعریف وچ
(\olref[sth][z][infinity-again]{sec} ویکھو)۔ ایہ اتفاق نہیں۔ فون نوئمان دی
ترتیبی عدداں دی تعریف قدرتی عدداں نوں ترتیبی عدد منّدی اے، تے لامتناہی
درجے دے ترتیبی عدداں دی وی تھاں رکھدی اے۔

ہمیشہ وانگ ہُن پُچھ سکدے آں کہ کی ایہ \emph{واقعی} ترتیبی عدد نیں؟
یا فون نوئمان نے بس کجھ سیٹ دتے نیں جیہناں نوں اسیں ترتیبی عدد
\emph{منّ} سکدے آں؟ ایس سوال اُتے جیہڑی گل بات ہو سکدی اے، اوہ
\olref[sfr][rel][ref]{sec}، \olref[sfr][arith][ref]{sec}،
\olref[sfr][infinite][dedekindsproof]{sec} تے \olref[sth][z][nat]{sec}
والی گل بات ورگی اے، سو اسیں ایس نوں نہیں کھچاں گے۔ اگے اسیں بس
’’ترتیبی عدد‘‘ نال ’’فون نوئمان دے ترتیبی عدد‘‘ آکھاں گے۔

\end{document}
